\documentclass[11pt,a4paper]{article}

\usepackage[T1]{fontenc}
\usepackage[utf8]{inputenc}
\usepackage{lmodern}
\usepackage{microtype}
\usepackage[margin=24mm]{geometry}
\usepackage{amsmath,amssymb,mathtools}
\usepackage{booktabs,longtable,array,tabularx}
\usepackage{enumitem}
\usepackage{xcolor}
\usepackage{listings}
\usepackage{hyperref}
\usepackage{fancyhdr}

\definecolor{nismoblue}{HTML}{164E63}
\definecolor{nismogreen}{HTML}{166534}
\definecolor{nismored}{HTML}{991B1B}
\definecolor{nismogray}{HTML}{F3F4F6}
\definecolor{codegray}{HTML}{4B5563}

\hypersetup{
  colorlinks=true,
  linkcolor=nismoblue,
  urlcolor=nismoblue,
  citecolor=nismoblue,
  pdftitle={NISMO: Mathematics, Implementation, API Adaptation, and Termination Rules},
  pdfauthor={NISMO contributors}
}

\lstdefinestyle{nismopython}{
  language=Python,
  basicstyle=\ttfamily\small,
  keywordstyle=\color{nismoblue}\bfseries,
  commentstyle=\color{nismogreen},
  stringstyle=\color{nismored},
  numbers=left,
  numberstyle=\scriptsize\color{codegray},
  numbersep=8pt,
  backgroundcolor=\color{nismogray},
  frame=single,
  rulecolor=\color{gray!35},
  breaklines=true,
  breakatwhitespace=false,
  showstringspaces=false,
  tabsize=4,
  columns=fullflexible,
  keepspaces=true
}
\lstset{style=nismopython}

\newcommand{\NISMO}{\textsc{NISMO}}
\newcommand{\MorphZ}{\textsc{MorphZ}}
\newcommand{\code}[1]{\texttt{\detokenize{#1}}}
\newcommand{\breakcode}[1]{\nolinkurl{#1}}
\newcommand{\R}{\mathbb{R}}
\newcommand{\dd}{\,\mathrm{d}}
\newcommand{\logZ}{\log Z}
\newcommand{\logPsi}{\log\Psi}
\newcommand{\important}[1]{%
  \begin{center}
  \fcolorbox{nismoblue}{nismogray}{%
    \begin{minipage}{0.92\linewidth}
    #1
    \end{minipage}}
  \end{center}}

\setlist{nosep,leftmargin=*}
\setlength{\parindent}{0pt}
\setlength{\parskip}{0.55em}
\setlength{\LTpre}{0.3em}
\setlength{\LTpost}{0.8em}
\renewcommand{\arraystretch}{1.18}

\pagestyle{fancy}
\fancyhf{}
\lhead{\NISMO{} mathematics and implementation}
\rhead{\thepage}
\setlength{\headheight}{14pt}

\title{\textbf{NISMO}\\
Mathematics, Implementation, API Adaptation,\\
and Termination Rules}
\author{Implementation guide for \code{nismo} version \code{0.1.0.dev3}}
\date{Code audit: repository commit \code{934195f}}

\begin{document}

\maketitle

\begin{abstract}
This document explains the mathematics implemented by the current
Morph-assisted nested importance sampling (\NISMO) code and connects each
equation to the corresponding stage of the Python execution path.  It also
specifies the public input contracts, shows which arguments must change for
common alternative user setups, describes the returned weighted posterior,
and gives the exact scientific, resource-limit, and exception termination
semantics.  The description is an implementation audit of the repository
version stated above; it is not a generic account of all algorithms that might
be called NISMO.
\end{abstract}

\tableofcontents
\newpage

\section{What the current code does}

\NISMO{} is a \emph{post-processing evidence estimator}.  The user first
provides posterior samples obtained by some other method.  \MorphZ{} fits a
fixed normalized density \(q(\theta)\) to those samples.  The sampler then
rewrites the Bayesian evidence integral as a nested-sampling problem under
\(q\), performs independent constrained rejection draws from that fixed
proposal, and constructs evidence and posterior weights using a deterministic
prior-volume trajectory.

At a high level, one run performs the following operations:

\begin{enumerate}
  \item validate the model, proposal, run limits, and stopping configuration;
  \item draw \(N\) new live points from the fixed proposal \(q\);
  \item evaluate the likelihood, normalized prior, proposal density, and
        transformed pseudo-likelihood at every point;
  \item repeatedly discard the live point with the smallest transformed
        pseudo-likelihood and replace it with an independent draw satisfying
        the new constraint;
  \item add the discarded point's deterministic quadrature contribution;
  \item calculate the current live contribution, information, error
        approximation, and stopping diagnostics;
  \item stop on a successful scientific rule or on a hard resource/rejection
        limit;
  \item add the final live-point correction and return an immutable result.
\end{enumerate}

\important{\textbf{Scope and central limitation.}
The current Phase~2 implementation does not discover a posterior from the
original prior.  It assumes that the input posterior samples already represent
every region of material posterior mass.  The fitted proposal is fixed and
non-defensive: it is not mixed with the prior and is not adapted during the
run.  Missing proposal support can therefore bias an evidence estimate.}

\section{Mathematical target and change of measure}
\label{sec:target}

\subsection{Bayesian evidence}

Let \(\theta\in\Theta\subseteq\R^d\) denote the parameters in the exact
coordinate system used by the code.  Let
\[
  L(\theta) = p(y\mid\theta)
  \qquad\text{and}\qquad
  \pi(\theta) = p(\theta)
\]
be the likelihood and a \emph{normalized} prior density with respect to the
declared base measure.  The evidence is
\begin{equation}
  Z = p(y) = \int_{\Theta} L(\theta)\pi(\theta)\dd\theta .
  \label{eq:evidence}
\end{equation}
All logarithms in the API, implementation, history, and this document are
natural logarithms.

The prior supplied to \code{log_prior_fn} must contain every normalization
constant.  For example, the normalized log density for a uniform prior on a
box with side lengths \(b_k-a_k\) is
\[
  \log\pi(\theta)=
  \begin{cases}
    -\sum_{k=1}^{d}\log(b_k-a_k),&
      a_k\leq\theta_k\leq b_k\ \text{for all }k,\\
    -\infty,&\text{otherwise}.
  \end{cases}
\]
Supplying only an ``inside/outside'' indicator would shift \(\log Z\) by the
missing log-volume.

For an absolute evidence calculation, \code{log_likelihood_fn} must likewise
represent the intended likelihood including any parameter-independent
normalization terms that would matter when comparing models.

\subsection{Fixed importance Morph and pseudo-likelihood}

The posterior training array is used to fit one normalized importance density
\(q_0(\theta)\).  In the standard path, \(q_0\) is a \MorphZ{}
\code{GroupKDE}; a custom object satisfying the same proposal protocol may
also be used.

The implementation defines
\begin{equation}
  \Psi_0(\theta)
  = \frac{L(\theta)\pi(\theta)}{q_0(\theta)},\qquad
  \log\Psi_0(\theta)
  = \log L(\theta)+\log\pi(\theta)-\log q_0(\theta).
  \label{eq:psi}
\end{equation}
Consequently,
\begin{equation}
  Z
  = \int_{\Theta}\Psi_0(\theta)q_0(\theta)\dd\theta.
  \label{eq:change-measure}
\end{equation}
Thus \(q_0\) acts as the nested sampler's pseudo-prior and \(\Psi_0\) acts as
its pseudo-likelihood.  This implementation fixes any tempering
parameter at \(\beta=1\); there is no public \code{beta} argument.

The default \code{fixed_morph} scheme also proposes from \(q_0\).  The
experimental \code{adaptive_morph} scheme refits a separate proposal from the
current live set every 25 iterations by default, while all density evaluation
and ordering continue to use \(q_0\) and \(\Psi_0\).  It directly accepts the
first proposal above the constraint without a Metropolis--Hastings correction.
It therefore does not in general preserve constrained \(q_0\), and its
quadrature is heuristic and may be biased.

The code computes Equation~\eqref{eq:psi} in
\code{BatchEvaluator.evaluate} after evaluating all three log densities.
It permits \(-\infty\) likelihood or prior values to represent zero density.
It rejects NaN and \(+\infty\).  If
\(\log L+\log\pi\) is finite but \(\log q_0=-\infty\), it raises
\code{ProposalSupportError}; the run does not return a result.

\subsection{Support and coordinate requirements}
\label{sec:support}

The identity in Equation~\eqref{eq:change-measure} is useful only if
\[
  q_0(\theta)>0
  \quad\text{wherever}\quad
  L(\theta)\pi(\theta)
\]
has material mass.  A finite numerator divided by zero proposal density is not
repairable by the current implementation.

The likelihood, prior, posterior training samples, and proposal density must
also use the same:

\begin{itemize}
  \item parameter order;
  \item units and transformations;
  \item Jacobian convention;
  \item dimensionality; and
  \item underlying density measure.
\end{itemize}

If a user changes from \(\theta\) to a transformed coordinate
\(\phi=g(\theta)\), all model and proposal inputs must change consistently.
In particular, the density in the new coordinates must include the relevant
Jacobian.  Merely transforming the samples without changing
\code{log_prior_fn} and \code{log_likelihood_fn} violates the implemented
identity.

\section{Nested-importance quadrature}
\label{sec:quadrature}

\subsection{Nested volume under the proposal}

For a pseudo-likelihood threshold \(\lambda\), define the remaining proposal
mass
\begin{equation}
  X(\lambda)
  = \int_{\Psi(\theta)>\lambda}q(\theta)\dd\theta.
\end{equation}
Under the usual nested-sampling ordering, the evidence can be viewed as
\[
  Z=\int_0^1 \Psi(X)\dd X .
\]

With \(N=\code{n_live}\) live points, the code does not draw a random shrinkage
factor for the quadrature.  It uses the deterministic expected log-volume
trajectory
\begin{equation}
  X_i = \exp(-i/N),\qquad
  \log X_i=-i/N,
  \label{eq:volume}
\end{equation}
where \(i=1,\ldots,K\) indexes completed replacements.  This choice is central:
changing a stopping policy does not change Equation~\eqref{eq:volume}.

\subsection{Dead-point contribution}

At iteration \(i\), let \(\Psi_i\) be the discarded minimum live
pseudo-likelihood.  The rectangular quadrature interval and contribution are
\begin{align}
  \Delta X_i &= X_{i-1}-X_i,\\
  w_i^{\mathrm{dead}} &= \Delta X_i\,\Psi_i,\\
  \log w_i^{\mathrm{dead}}
  &= \log\!\left(X_{i-1}-X_i\right)+\log\Psi_i.
  \label{eq:dead-weight}
\end{align}
\code{dead_log_contribution} evaluates these quantities in log space using a
stable log-difference.

\subsection{Final live correction}

After \(K\) completed replacements, the remaining proposal volume is
\(X_K=\exp(-K/N)\).  The code assigns equal volume \(X_K/N\) to each current
live point:
\begin{align}
  w_j^{\mathrm{live}}
    &= \frac{X_K}{N}\Psi_j^{\mathrm{live}},\\
  \log w_j^{\mathrm{live}}
    &= \log X_K-\log N+\log\Psi_j^{\mathrm{live}},
    \qquad j=1,\ldots,N.
  \label{eq:live-weight}
\end{align}
This correction is applied on every returned result, including a hard-stop
partial result.

\subsection{Evidence, posterior weights, and information}

Let \(a\) run over the \(K\) dead contributions followed by the \(N\) final
live contributions.  Finalization computes
\begin{align}
  \log\widehat Z
    &= \operatorname{logsumexp}_a(\log w_a),\\
  \log\widetilde w_a
    &= \log w_a-\log\widehat Z,\\
  \widetilde w_a
    &= \exp(\log\widetilde w_a).
  \label{eq:posterior-weights}
\end{align}
These normalized quadrature contributions are the returned posterior weights.
The ordering is always all dead points first and all final live points second.
The primary posterior representation is therefore
\[
  \bigl(\code{result.all_points},
        \code{result.posterior_weights}\bigr).
\]

The discrete information estimate is
\begin{equation}
  \widehat H
  = \sum_a \widetilde w_a
      \left(\log\Psi_a-\log\widehat Z\right),
  \label{eq:information}
\end{equation}
with tiny negative floating-point roundoff clipped to zero.  The returned
theoretical nested-sampling error approximation is
\begin{equation}
  \code{logzerr}
  = \sqrt{\widehat H/N}.
  \label{eq:logzerr}
\end{equation}
It is an approximation rather than a calibration guarantee.  In particular,
it does not include Morph fitting uncertainty, missing modes or support,
stochastic uncertainty that deterministic shrinkage omits, or invalid
constrained draws.

\subsection{Equal-weight resampling}

\code{NISMOResult.resample_equal} uses randomized systematic resampling.  For a
requested count \(M\), it draws one \(U\sim\operatorname{Uniform}(0,1)\), forms
\[
  r_m=\frac{m+U}{M},\qquad m=0,\ldots,M-1,
\]
maps these positions through the cumulative posterior weights, and shuffles
the selected indices.  Repeated points are expected.  This adds resampling
noise, so weighted summaries should be preferred whenever the downstream tool
accepts weights.

\section{Exact execution path}
\label{sec:execution}

This section follows \code{NISMOSampler.run} in execution order.

\subsection{Step 0: resolve and validate configuration}

\code{NISMOConfig} is immutable and validates all settings before sampling.
If neither \code{dlogz} nor \code{stopping} is supplied, the configuration is
resolved to \code{dlogz=1e-3}.  A \code{dlogz} value is internally
represented as a one-criterion \code{StoppingPolicy} using
\code{remaining_dlogz}.  Supplying both interfaces is an error.

The run then creates the progress reporter, records the initial random-number
generator state and start time, derives a wall-time deadline if requested, and
constructs a counting \code{BatchEvaluator}.

\subsection{Step 1: initialize the live set}

The sampler requests exactly \(N\) \emph{new} points from
\code{proposal.sample}.  It never installs the Morph training rows directly as
live points.  The returned array must be finite and have shape \((N,d)\).

For each live point, the evaluator caches
\[
  \theta,\quad \log L,\quad \log\pi,\quad \log q,\quad \log\Psi.
\]
It also draws a uniform tie breaker \(u_j\in[0,1)\) for every live point,
regardless of which tie policy is selected.  The initial evaluation consumes
\(N\) likelihood calls and \(N\) prior calls.  Initialization happens before
the sampler checks run-time hard limits inside the main loop.

\subsection{Step 2: check hard limits at the top of the loop}

Before selecting a point for the next replacement, the sampler checks, in this
order:

\begin{enumerate}
  \item \code{niter >= max_iterations};
  \item the run-wide likelihood count has reached
        \code{max_likelihood_calls};
  \item the monotonic clock has reached \code{max_wall_time}.
\end{enumerate}

The first true condition determines the hard-stop reason.  These checks occur
before each attempted replacement, not midway through evidence arithmetic.

\subsection{Step 3: choose the worst live point}

For the default \code{tie_policy="strict"}, the code selects
\[
  j_\star=\arg\min_j\log\Psi_j
\]
using \code{numpy.argmin}.  This is appropriate for ordinary continuous
pseudo-likelihoods, for which exact ties have probability zero.

For \code{tie_policy="randomized_plateau"}, it orders the pairs
\[
  (\log\Psi_j,u_j)
\]
lexicographically and removes the smallest pair.  The auxiliary uniform value
provides a continuous ordering within an exact pseudo-likelihood plateau.

\subsection{Step 4: draw an independent constrained replacement}
\label{sec:constrained}

The code repeatedly draws independent batches from \(q\).  For strict
ordering, a candidate is valid when
\[
  \log\Psi_{\mathrm{candidate}}
  > \log\Psi_{j_\star}.
\]
For randomized plateau ordering, it is valid when
\[
  \log\Psi_{\mathrm{candidate}}
    > \log\Psi_{j_\star}
  \quad\text{or}\quad
  \left(
    \log\Psi_{\mathrm{candidate}}=\log\Psi_{j_\star}
    \quad\text{and}\quad
    u_{\mathrm{candidate}}>u_{j_\star}
  \right).
\]

Valid candidates are scanned in generation order and the first is accepted.
The code never selects the largest candidate in a batch.  Therefore, under
the stated independence and density assumptions, strict rejection targets
\[
  q\!\left(\theta\mid
    \log\Psi(\theta)>\log\Psi_{j_\star}\right).
\]
The randomized rule targets the analogous constrained distribution on the
augmented \((\theta,u)\) space.

An implementation detail relevant to cost accounting is that the entire
generated batch is evaluated before the first valid row is selected.
Consequently, all rows in that batch count as proposals and likelihood calls,
even when an early row is accepted.

\subsection{Step 5: record the dead point and replace it}

After a successful constrained draw, iteration \(i=K+1\) is completed:

\begin{enumerate}
  \item compute \(\log X_i\), \(\log\Delta X_i\), and the dead log
        contribution from Equation~\eqref{eq:dead-weight};
  \item append the discarded point and all its cached density values;
  \item replace its live-set row with the accepted point and its tie breaker;
  \item update the cumulative dead evidence in log space;
  \item increment \code{niter}.
\end{enumerate}

Because the replacement satisfies the selected ordering rule, stored dead
thresholds should be monotone.

\subsection{Step 6: estimate the current live contribution}

Using the post-replacement live set and \(X_i\), the live evidence estimate is
\begin{equation}
  \log \widehat Z_{\mathrm{live},i}
  = \log X_i+
    \operatorname{logsumexp}_{j=1}^{N}
       (\log\Psi_j^{\mathrm{live}})
    -\log N.
  \label{eq:live-logz}
\end{equation}
The current total is
\[
  \log\widehat Z_{\mathrm{total},i}
  =\operatorname{logaddexp}
     \left(\log\widehat Z_{\mathrm{dead},i},
           \log\widehat Z_{\mathrm{live},i}\right).
\]
The code calculates the information and
\(\sqrt{\widehat H/N}\) at the same point in the iteration.

\subsection{Step 7: calculate stopping diagnostics and update the streak}

Every stopping metric in Section~\ref{sec:stopping-metrics} is calculated after
the successful replacement.  Enabled criteria are evaluated using inclusive
comparisons (\(\leq\) or \(\geq\)).  The configured \code{all}/\code{any}
combination is gated by \code{min_iterations}.  A passing combination
increments the streak; any failure, including being below
\code{min_iterations}, resets it to zero.

All history arrays and an optional progress callback are updated before a
scientific stop is installed.  Thus the final successful iteration appears in
\code{result.history}.

\subsection{Step 8: finalize on every returned path}

When the loop ends with a scientific or hard-stop reason, the code:

\begin{enumerate}
  \item sets \(\log X_K=-K/N\);
  \item combines all dead weights with every live weight from
        Equation~\eqref{eq:live-weight};
  \item recomputes final \(\log Z\), \(H\), \code{logzerr}, and normalized
        posterior weights;
  \item freezes result and history arrays as read-only;
  \item stores the complete configuration, random-generator states, proposal
        metadata representation, counts, warnings, and termination status.
\end{enumerate}

The result object validates that posterior weights normalize, thresholds are
monotone, and the stored evidence can be reconstructed from the stored arrays.

\section{Public API and what changes for a different setup}
\label{sec:api}

\subsection{Baseline end-to-end example}

\begin{lstlisting}
import numpy as np

from nismo import (
    CallableModel,
    NISMOSampler,
    MorphProposal,
    StoppingCriterionConfig,
    StoppingPolicy,
)

def log_likelihood(theta):
    # theta has shape (n, 2); return shape (n,)
    return -np.log(2.0 * np.pi) - 0.5 * np.sum(theta**2, axis=1)

def normalized_log_prior(theta):
    inside = np.all(np.abs(theta) <= 5.0, axis=1)
    return np.where(inside, -2.0 * np.log(10.0), -np.inf)

model = CallableModel(
    ndim=2,
    parameter_names=("x", "y"),
    log_likelihood_fn=log_likelihood,
    log_prior_fn=normalized_log_prior,
)

# Rows must use the same (x, y) coordinates as the model.
posterior_samples = np.load("posterior_samples.npy")

importance_morph = MorphProposal.fit(
    posterior_samples,
    morph_type="2_group",
    param_names=model.parameter_names,
    kde_bw="silverman",
)

sampler = NISMOSampler(
    model=model,
    importance_morph=importance_morph,
    proposal_scheme="fixed_morph",
    proposal_update_interval=25,
    n_live=500,
    rng=42,
    proposal_batch_size=64,
    tie_policy="strict",
)

result = sampler.run(
    dlogz=1.0e-3,
    max_iterations=20_000,
    max_proposals_per_replacement=200_000,
    max_likelihood_calls=None,
    max_wall_time=None,
    progress=True,
)

if not result.success:
    raise RuntimeError(result.termination_reason)

posterior_mean = np.average(
    result.all_points,
    axis=0,
    weights=result.posterior_weights,
)
\end{lstlisting}

\subsection{Model arguments}

\begin{longtable}{>{\raggedright\arraybackslash}p{0.20\linewidth}
                  >{\raggedright\arraybackslash}p{0.23\linewidth}
                  >{\raggedright\arraybackslash}p{0.49\linewidth}}
\toprule
\textbf{Argument} & \textbf{Contract} & \textbf{Change it when}\\
\midrule
\endfirsthead
\toprule
\textbf{Argument} & \textbf{Contract} & \textbf{Change it when}\\
\midrule
\endhead
\code{ndim} &
Positive integer \(d\). &
The parameter dimension changes.  Then sample arrays must have \(d\) columns,
\code{parameter_names} must have \(d\) unique entries, and
\code{proposal.ndim} must equal \(d\).\\
\code{parameter_names} &
Unique names in column order. &
The variables or their order change.  The same order must be used by the
posterior array and Morph grouping definitions.\\
\breakcode{log_likelihood_fn} &
Returns \((n,)\) log-likelihood values. &
The data model changes.  It receives \((n,d)\) when vectorized or a single
\((d,)\) row when \code{vectorized=False}.\\
\code{log_prior_fn} &
Returns a normalized \((n,)\) log-prior. &
Prior bounds, distribution, coordinates, or dimension change.  Include all
normalization and Jacobian terms; use \(-\infty\) for zero density.\\
\code{vectorized} &
Default \code{True}. &
Set to \code{False} only when \emph{both} supplied functions accept individual
rows rather than batches.  The adapter does not catch exceptions to guess
vectorization.\\
\bottomrule
\end{longtable}

For scalar user functions, the setup becomes:

\begin{lstlisting}
def scalar_log_likelihood(theta_row):
    return -0.5 * theta_row[0] ** 2

def scalar_log_prior(theta_row):
    return -0.5 * theta_row[0] ** 2 - 0.5 * np.log(2.0 * np.pi)

model = CallableModel(
    ndim=1,
    parameter_names=("x",),
    log_likelihood_fn=scalar_log_likelihood,
    log_prior_fn=scalar_log_prior,
    vectorized=False,
)
\end{lstlisting}

\subsection{Morph fitting arguments}
\label{sec:morph-args}

\begin{longtable}{>{\raggedright\arraybackslash}p{0.20\linewidth}
                  >{\raggedright\arraybackslash}p{0.27\linewidth}
                  >{\raggedright\arraybackslash}p{0.45\linewidth}}
\toprule
\textbf{Argument} & \textbf{Current contract/default} &
\textbf{Effect and setup change}\\
\midrule
\endfirsthead
\toprule
\textbf{Argument} & \textbf{Current contract/default} &
\textbf{Effect and setup change}\\
\midrule
\endhead
\breakcode{posterior_samples} &
Finite array \((n_{\rm train},d)\), at least two rows. &
Replace with representative samples for the new posterior and coordinate
system.  Training data are copied.  There is no weights argument: weighted
external samples must be converted to an appropriate unweighted training
sample before this call.\\
\code{morph_type} &
\code{None} or \code{"<k>_group"}, with \(2\leq k\leq d\). &
Use, for example, \code{"2_group"} to compute \(k\)-order total correlations
and let \MorphZ{} greedily choose disjoint groups.  The literal
\code{"n_group"} is invalid.\\
\code{group_file} &
\code{None} or path to Morph-compatible JSON. &
Use when a grouping definition was computed previously.  NISMO reads it and
passes its contents in memory.\\
\code{groups} &
\code{None} or in-memory Morph grouping definition. &
Use for explicit groups.  \code{groups=[]} requests independent
one-dimensional KDE components.\\
Grouping choice &
At most one of \code{morph_type}, \code{group_file}, and \code{groups}. &
If all are omitted, the implementation also uses independent one-dimensional
components.  Correlated targets generally require a grouping choice that can
represent those correlations.\\
\code{param_names} &
Defaults to \code{param_0}, \code{param_1}, \ldots. &
Normally pass \code{model.parameter_names}.  Names must be unique and match
sample columns and group entries.\\
\code{kde_bw} &
\code{"silverman"}; string, float, or per-parameter dictionary. &
Change the KDE smoothing rule for the new training set.  It is forwarded
unchanged to \code{morphZ.GroupKDE}; it changes \(q\), support behavior, and
rejection efficiency.\\
\code{min_tc} &
\code{None}. &
Forwarded to \MorphZ{} as the total-correlation selection threshold.\\
\code{top_k_greedy} &
\code{1}, positive integer. &
Forwarded to the greedy group selection.\\
\bottomrule
\end{longtable}

The three grouping patterns are:

\begin{lstlisting}
# Automatic second-order total-correlation workflow
proposal = MorphProposal.fit(
    samples,
    morph_type="2_group",
    param_names=model.parameter_names,
)

# Precomputed JSON groups
proposal = MorphProposal.fit(
    samples,
    group_file="groups.json",
    param_names=model.parameter_names,
)

# Explicit independent one-dimensional components
proposal = MorphProposal.fit(
    samples,
    groups=[],
    param_names=model.parameter_names,
)
\end{lstlisting}

\code{NISMOSampler.from_posterior_samples} is a convenience path.  Every item in
\code{morph_config} is passed to \code{MorphProposal.fit}, while parameter
names are taken from the model; do not duplicate \code{param_names} inside
\code{morph_config}:

\begin{lstlisting}
sampler = NISMOSampler.from_posterior_samples(
    model=model,
    posterior_samples=samples,
    morph_config={
        "morph_type": "2_group",
        "kde_bw": "silverman",
    },
    n_live=500,
    rng=42,
)
\end{lstlisting}

\subsection{Using a non-Morph proposal}

The sampler is typed against a small normalized proposal protocol.  A custom
proposal must expose \code{ndim}, draw independent finite \((n,d)\) samples,
and return normalized \((n,)\) log densities:

\begin{lstlisting}
class UniformBoxProposal:
    ndim = 2

    def sample(self, n, rng):
        return rng.uniform(-5.0, 5.0, size=(n, self.ndim))

    def log_prob(self, theta):
        inside = np.all(np.abs(theta) <= 5.0, axis=1)
        return np.where(inside, -2.0 * np.log(10.0), -np.inf)

importance_morph = UniformBoxProposal()
sampler = NISMOSampler(
    model=model,
    importance_morph=importance_morph,
    n_live=500,
    rng=42,
)
\end{lstlisting}

In this setup, omit every \MorphZ{} fitting argument.  The change-of-measure
identity still requires the custom \code{log_prob} to be the normalized density
of the distribution actually sampled by \code{sample}.  Correlation among
coordinates within one multivariate draw is allowed.  Dependence between
proposal rows or across successive draws, or state-dependent sampling,
violates the current independent rejection argument.

\subsection{Sampler-construction arguments}

\begin{longtable}{>{\raggedright\arraybackslash}p{0.22\linewidth}
                  >{\raggedright\arraybackslash}p{0.25\linewidth}
                  >{\raggedright\arraybackslash}p{0.45\linewidth}}
\toprule
\textbf{Argument} & \textbf{Validation/default} & \textbf{Practical effect}\\
\midrule
\endfirsthead
\toprule
\textbf{Argument} & \textbf{Validation/default} & \textbf{Practical effect}\\
\midrule
\endhead
\code{model} &
Implements the model protocol. &
Defines \(L\), normalized \(\pi\), dimension, and parameter order.\\
\code{proposal} &
Normalized proposal with matching \code{ndim}. &
Defines the fixed \(q\), initial live distribution, and constrained rejection
source.\\
\code{n_live} &
Required integer \(N\geq2\). &
Changes the shrinkage schedule \(X_i=e^{-i/N}\), initial evaluation cost,
live-ESS range, resolution, theoretical error scale, and usually runtime.
Larger is normally more expensive and more stable.\\
\code{rng} &
Required Python integer seed or \breakcode{numpy.random.Generator}. &
Controls proposal resampling and tie breakers.  A supplied generator is
consumed in place.  Recreate the sampler with the same seed and configuration
for reproducibility.\\
\breakcode{proposal_batch_size} &
Default \(64\), positive integer. &
Changes vectorized rejection throughput.  A full batch is evaluated and
counted.  It should fit model memory.  Changing it can change the random stream
and realized run even with the same starting seed.\\
\code{tie_policy} &
\code{"strict"} by default; alternatively
\breakcode{"randomized_plateau"}. &
Use strict ordering for continuous \(\log\Psi\).  Use randomized plateau
ordering when exact ties have nonzero probability, such as discrete or
piecewise-constant targets.\\
\bottomrule
\end{longtable}

\subsection{Run arguments}

\begin{longtable}{>{\raggedright\arraybackslash}p{0.24\linewidth}
                  >{\raggedright\arraybackslash}p{0.24\linewidth}
                  >{\raggedright\arraybackslash}p{0.44\linewidth}}
\toprule
\textbf{Argument} & \textbf{Validation/default} & \textbf{When to change it}\\
\midrule
\endfirsthead
\toprule
\textbf{Argument} & \textbf{Validation/default} & \textbf{When to change it}\\
\midrule
\endhead
\code{dlogz} &
\code{None}; if both stopping inputs are omitted, resolves to \(10^{-3}\).
An explicit value must be positive and finite. &
Use for the single \code{remaining_dlogz} scientific rule.  It bounds the
estimated change in accumulated log evidence from adding the mean-live
remainder.  Smaller values generally run longer.\\
\code{stopping} &
\code{None} or a \code{StoppingPolicy}; mutually exclusive with \code{dlogz}. &
Use to combine live uncertainty, remaining mass, stability, ESS, or
theoretical error.\\
\code{max_iterations} &
Default \(10{,}000\), positive integer. &
Raise when a scientifically adequate run needs more replacements.  Hitting
the limit is failure, not convergence.\\
\breakcode{max_proposals_per_replacement} &
Default \(100{,}000\), positive integer. &
Raise when late constrained rejection is valid but rare.  Hitting it produces
\code{constrained_sampling_exhausted}, or \code{plateau_stall} in the
specific strict-tie case.\\
\breakcode{max_likelihood_calls} &
Default \code{None}; otherwise integer at least \(N\). &
Set a run-wide model-evaluation budget.  The count includes the initial
\(N\) live points and every evaluated proposal row.\\
\code{max_wall_time} &
Default \code{None}; otherwise positive seconds. &
Set a monotonic wall-clock budget.  It is checked between batches and at the
top of iterations, not by interrupting a user density function.\\
\code{progress} &
\code{False}, \code{True}, or callable. &
Use \code{True} for the optional \code{tqdm} display, or pass a callback for
structured snapshots.  It does not change quadrature.\\
\bottomrule
\end{longtable}

\subsection{Setup-change checklist}
\label{sec:setup-checklist}

\begin{longtable}{>{\raggedright\arraybackslash}p{0.29\linewidth}
                  >{\raggedright\arraybackslash}p{0.63\linewidth}}
\toprule
\textbf{New setup} & \textbf{Arguments and contracts that must change}\\
\midrule
\endfirsthead
\toprule
\textbf{New setup} & \textbf{Arguments and contracts that must change}\\
\midrule
\endhead
Different dimension or variables &
Change \code{ndim}, \code{parameter_names}, both density functions, posterior
sample columns, and proposal dimension.  Refit the proposal.  Reconsider
\code{morph_type} because \(k\leq d\).\\
Different prior &
Change \code{log_prior_fn}, including normalization and any Jacobian.  Ensure
the posterior training data and \(q\) cover the new posterior.\\
Different likelihood or data &
Change \code{log_likelihood_fn}; regenerate representative posterior training
samples and refit \(q\).\\
Scalar rather than batched functions &
Set \code{vectorized=False} and make both callables accept one \((d,)\) row.
Expect lower throughput.\\
Correlated posterior &
Use an appropriate \code{morph_type}, \code{group_file}, or nonempty
\code{groups}; do not pass more than one.  Tune \code{kde_bw} and verify
support and rejection efficiency.\\
Independent posterior components &
Use \code{groups=[]}, or omit all grouping arguments (the same independent
default).\\
Weighted posterior input &
There is no Morph weights API in this adapter.  Produce a representative
unweighted training array externally before
\code{MorphProposal.fit}.\\
Exact likelihood/pseudo-likelihood plateaus &
Set \code{tie_policy="randomized_plateau"}.  Strict ordering can exhaust its
proposal budget because equal values do not satisfy \(>\).\\
Custom normalized proposal &
Implement \code{ndim}, \code{sample(n,rng)}, and \code{log_prob(theta)};
construct \code{NISMOSampler} directly and omit Morph fitting arguments.\\
More stringent default completion &
Decrease \code{dlogz}; also raise hard budgets so they do not stop the run
first.  Validate by repeated runs rather than interpreting \code{dlogz} as an
absolute evidence error.\\
Hybrid completion rule &
Omit \code{dlogz} and pass a \code{StoppingPolicy}.  Choose tolerances,
\code{mode}, \code{consecutive}, \code{min_iterations}, and
\code{stability_window} based on repeated benchmarks for the target family.\\
Expensive vectorized model &
Tune \code{proposal_batch_size} to the model's throughput and memory; use
\code{max_likelihood_calls} or \code{max_wall_time} as safeguards.\\
Reproducible comparison &
Use explicit seeds, keep \code{proposal_batch_size} and all policies fixed,
record \code{result.config}, and compare complete independent runs.\\
\bottomrule
\end{longtable}

\section{Current scientific stopping rules}
\label{sec:stopping-metrics}

\subsection{Remaining log-evidence increment}

The estimated increment from adding the mean-live remainder to the accumulated
dead evidence is
\begin{equation}
  \Delta\log Z_{\mathrm{rem}}
  = \log\left(\widehat Z_{\mathrm{dead}}+\widehat Z_{\mathrm{live}}\right)
    -\log\widehat Z_{\mathrm{dead}}
  = \log\left(
      1+\frac{\widehat Z_{\mathrm{live}}}{\widehat Z_{\mathrm{dead}}}
    \right).
  \label{eq:remaining-dlogz}
\end{equation}
The criterion
\[
  \code{remaining_dlogz}\leq\tau
\]
passes at equality and requires a positive finite tolerance; \(\tau\) need not
be less than one.  Before finite positive dead evidence has accumulated, the
diagnostic is \(+\infty\) and cannot pass.

The call
\begin{lstlisting}
result = sampler.run(dlogz=1.0e-3)
\end{lstlisting}
is exactly a one-criterion \code{remaining_dlogz} policy with
\(\tau=\code{dlogz}\), \code{mode="all"}, one required consecutive pass, and
no minimum iteration gate.  Omitting both \code{dlogz} and \code{stopping}
selects this same \(10^{-3}\) behavior.

\subsection{Remaining live-evidence fraction}

The live fraction is
\begin{equation}
  f_{\mathrm{live}}
  = \frac{\widehat Z_{\mathrm{live}}}
         {\widehat Z_{\mathrm{total}}}
  = \exp\left(
      \log\widehat Z_{\mathrm{live}}
      -\log\widehat Z_{\mathrm{total}}
    \right).
  \label{eq:remaining}
\end{equation}
The criterion
\[
  \code{remaining_fraction}\leq\tau
\]
passes at equality.  Its tolerance must satisfy \(0<\tau<1\).  This remains an
independently selectable explicit criterion.

The exact relationship between the two diagnostics is
\begin{equation}
  \Delta\log Z_{\mathrm{rem}}
  = -\log(1-f_{\mathrm{live}}).
\end{equation}
They are similar only for a small remainder and must not be conflated.

\subsection{Live effective sample size}

For \(a_j=\Psi_j^{\mathrm{live}}\), the code computes the Kish effective sample
size in log space:
\begin{equation}
  N_{\mathrm{eff,live}}
  = \frac{\left(\sum_{j=1}^{N}a_j\right)^2}
         {\sum_{j=1}^{N}a_j^2}.
  \label{eq:live-ess}
\end{equation}
The value is clipped to \([1,N]\).  Equal finite live contributions give
\(N_{\mathrm{eff,live}}=N\).  When every live contribution is zero
(\(\log\Psi_j=-\infty\)), the implementation also defines the live ESS as
\(N\), the remaining fraction as zero, the live-mean RSE as zero, and the live
log-evidence error as zero, provided the surrounding evidence state is
numerically consistent.

The selectable criterion is
\[
  \code{live_ess}\geq\tau,
\]
with \(1\leq\tau\leq N\).  It is the only currently implemented criterion that
uses a greater-than-or-equal comparison.

\subsection{Live-mean relative standard error}

The code derives
\begin{equation}
  \operatorname{RSE}_{\mathrm{live}}
  =
  \sqrt{
    \max\left(
      \frac{N/N_{\mathrm{eff,live}}-1}{N-1},
      0
    \right)
  }.
  \label{eq:rse}
\end{equation}
This is always recorded as \code{live_mean_rse}, but it is not a separately
selectable criterion in the current public policy names.

\subsection{Finite-live contribution to log-evidence error}

The first-order delta-method diagnostic is
\begin{equation}
  \sigma_{\log Z,\mathrm{live}}
  \approx
  f_{\mathrm{live}}\,
  \operatorname{RSE}_{\mathrm{live}}.
  \label{eq:live-logz-error}
\end{equation}
The criterion
\[
  \code{live_logz_error}\leq\tau
\]
requires \(\tau>0\).  It estimates only uncertainty transmitted from
representing the remaining integral by the finite live-set mean.  It excludes
proposal fitting, support, mode coverage, shrinkage, and constrained-draw
validity.  Equal live contributions produce zero empirical live-mean
uncertainty, so this rule is unsafe on its own.  The recommended policy pairs
it with a \code{remaining_dlogz} completion guard.

\subsection{Recent evidence stability}

For a configured window \(W=\code{stability_window}\), the diagnostic is
\begin{equation}
  S_i
  = \max_{r=i-W+1,\ldots,i}\log\widehat Z_r
    -\min_{r=i-W+1,\ldots,i}\log\widehat Z_r.
  \label{eq:stability}
\end{equation}
The value is NaN and the criterion cannot pass until exactly \(W\) completed
estimates are available.  The selectable condition
\[
  \code{logz_stability}\leq\tau
\]
requires \(\tau>0\) and \(W\geq2\).  Stability is supporting evidence only: a
biased or mode-missing estimate can be stable.

\subsection{Theoretical nested-sampling error}

The condition
\[
  \code{logzerr}\leq\tau
\]
uses Equation~\eqref{eq:logzerr} and requires \(\tau>0\).  It inherits all
limitations of the deterministic-shrinkage information approximation and is
not a complete calibration guarantee.

\subsection{Criterion summary}

\begin{center}
\begin{tabular}{lll}
\toprule
\textbf{Name} & \textbf{Passing comparison} & \textbf{Tolerance validation}\\
\midrule
\code{remaining_fraction} & \(f_{\mathrm{live}}\leq\tau\) & \(0<\tau<1\)\\
\code{remaining_dlogz} & \(\Delta\log Z_{\mathrm{rem}}\leq\tau\) & \(\tau>0\)\\
\code{live_logz_error} & \(\sigma_{\log Z,\mathrm{live}}\leq\tau\) & \(\tau>0\)\\
\code{logz_stability} & \(S_i\leq\tau\) & \(\tau>0\)\\
\code{live_ess} & \(N_{\mathrm{eff,live}}\geq\tau\) & \(1\leq\tau\leq N\)\\
\code{logzerr} & \(\sqrt{\widehat H/N}\leq\tau\) & \(\tau>0\)\\
\bottomrule
\end{tabular}
\end{center}

Criterion names must be unique within one policy.  All tolerances must be
finite real numbers and Boolean values are rejected.

\subsection{Combining, gating, and persisting criteria}
\label{sec:policy}

For enabled Boolean criterion results \(c_1,\ldots,c_m\):
\[
  c_{\mathrm{criteria}}=
  \begin{cases}
    \bigwedge_{r=1}^{m}c_r,&\code{mode="all"},\\
    \bigvee_{r=1}^{m}c_r,&\code{mode="any"}.
  \end{cases}
\]
The combined condition at completed iteration \(i\) is
\[
  c_i = [i\geq\code{min_iterations}]
        \wedge c_{\mathrm{criteria}}.
\]
The run-local streak is updated exactly as
\[
  s_i =
  \begin{cases}
    s_{i-1}+1,&c_i\text{ is true},\\
    0,&c_i\text{ is false}.
  \end{cases}
\]
Scientific stopping occurs when
\[
  s_i\geq\code{consecutive}.
\]
\code{consecutive} must be a positive integer,
\code{min_iterations} must be an integer at least zero, and
\code{stability_window} must be an integer at least two.  Even when
\code{min_iterations=0}, the earliest scientific stop is after the first
completed replacement because policies are not evaluated on the initial live
set.

A current experimental hybrid example is:

\begin{lstlisting}
n_live = sampler.n_live
stopping = StoppingPolicy(
    criteria=(
        StoppingCriterionConfig("remaining_dlogz", 1.0e-2),
        StoppingCriterionConfig("live_logz_error", 2.0e-3),
        StoppingCriterionConfig("logz_stability", 5.0e-3),
    ),
    mode="all",
    consecutive=3,
    min_iterations=n_live,
    stability_window=n_live,
)

result = sampler.run(
    stopping=stopping,       # do not also pass dlogz
    max_iterations=20_000,
)
\end{lstlisting}

These threshold values are illustrative calibration choices, not universal
constants.  They require repeated-run calibration using independent seeds,
known benchmarks where possible, proposal diagnostics, and a live count
appropriate to the new target.

\section{Exact termination semantics}
\label{sec:termination}

\subsection{Successful scientific termination}

There are exactly two termination reasons for which the result sets
\code{success=True}:

\begin{longtable}{>{\raggedright\arraybackslash}p{0.27\linewidth}
                  >{\raggedright\arraybackslash}p{0.65\linewidth}}
\toprule
\textbf{\code{termination_reason}} & \textbf{Exact meaning}\\
\midrule
\endfirsthead
\toprule
\textbf{\code{termination_reason}} & \textbf{Exact meaning}\\
\midrule
\endhead
\code{remaining_evidence} &
The run used the default/legacy-compatible \code{dlogz} interface, and after a
completed replacement its remaining-log-evidence-increment policy reached the
required streak.\\
\code{stopping_criteria} &
The run used an explicit \code{StoppingPolicy}, and after a completed
replacement its combined, gated, consecutive rule passed.\\
\bottomrule
\end{longtable}

Scientific stopping is checked only after a replacement is completed and
recorded.  There is no scientific stopping check on the initial live set
before the first replacement.

\subsection{Returned hard-stop failures}

The following reasons return an immutable, internally consistent
\code{NISMOResult} with the current final-live correction, but set
\code{success=False}:

\begin{longtable}{>{\raggedright\arraybackslash}p{0.32\linewidth}
                  >{\raggedright\arraybackslash}p{0.60\linewidth}}
\toprule
\textbf{\code{termination_reason}} & \textbf{Trigger}\\
\midrule
\endfirsthead
\toprule
\textbf{\code{termination_reason}} & \textbf{Trigger}\\
\midrule
\endhead
\code{max_iterations} &
At the top of an iteration, completed replacements satisfy
\code{niter >= max_iterations}.\\
\code{max_likelihood_calls} &
The run-wide count is exhausted either at the top of the loop or before a new
constrained batch.  The next batch size is capped by the remaining budget, so
the configured count is not intentionally exceeded.\\
\code{max_wall_time} &
The monotonic deadline is reached at the top of the loop or between constrained
proposal batches.  A running user function is not interrupted.\\
\breakcode{constrained_sampling_exhausted} &
One replacement consumes \code{max_proposals_per_replacement} evaluated
proposal rows without a valid constrained candidate.\\
\code{plateau_stall} &
The constrained proposal limit is exhausted under strict ordering and more
than one current live point equals the minimum threshold.  This is a
diagnostic relabeling of that exhausted attempt, indicating an exact-tie
plateau.\\
\bottomrule
\end{longtable}

A hard-stop evidence value is a \emph{partial-run estimate}.  It includes all
completed dead contributions and the current mean-live remainder, but a hard
limit is not evidence of convergence.  Code consuming results should always
check both fields:

\begin{lstlisting}
if not result.success:
    print("NISMO did not converge scientifically:", result.termination_reason)
\end{lstlisting}

\subsection{Ordering and simultaneous conditions}

At the top of a new loop pass, hard checks have the priority order
\code{max_iterations}, then \code{max_likelihood_calls}, then
\code{max_wall_time}.  Inside constrained rejection, wall time is checked
before the remaining likelihood budget for each new batch.

After a successful replacement, the scientific policy is evaluated
immediately.  If it passes, its scientific reason is installed before the next
top-of-loop hard checks.  For example, if an iteration both reaches
\code{max_iterations} and satisfies the scientific rule, the scientific
termination wins because the iteration-limit check would occur only on the
next pass.

\subsection{Raised errors: no result is returned}
\label{sec:exceptions}

Some failures mean that no statistically coherent result can be formed.  They
raise an exception rather than returning \code{success=False}:

\begin{longtable}{>{\raggedright\arraybackslash}p{0.32\linewidth}
                  >{\raggedright\arraybackslash}p{0.60\linewidth}}
\toprule
\textbf{Exception class/category} & \textbf{Examples}\\
\midrule
\endfirsthead
\toprule
\textbf{Exception class/category} & \textbf{Examples}\\
\midrule
\endhead
\code{ConfigurationError} &
Invalid live count, limits, tie policy, stopping name/tolerance, duplicate
criteria, or simultaneous \code{dlogz} and \code{stopping}.\\
\code{InvalidModelOutput} &
Wrong output shape, NaN, \(+\infty\), nonfinite input point, or invalid
computed \(\log\Psi\).\\
\code{InvalidProposalOutput} &
Wrong sample/density shape, nonfinite samples, NaN, or \(+\infty\) proposal
log density.\\
\code{ProposalSupportError} &
Finite \(\log L+\log\pi\) with \(\log q=-\infty\).\\
\code{NumericalInvariantError} &
An impossible evidence/metric state, nonfinite final evidence, materially
negative information, or posterior weights that do not normalize.\\
\breakcode{MissingOptionalDependency} &
\MorphZ{} is unavailable when fitting a Morph proposal, or \code{tqdm} is
unavailable when \code{progress=True}.\\
User exception &
An exception raised by a supplied likelihood, prior, proposal, progress
callback, or backend is not guessed away by the adapter.\\
\bottomrule
\end{longtable}

\section{Result interpretation and diagnostics}

\subsection{Core scalars and status}

\begin{itemize}
  \item \code{logz}, \code{information}, and \code{logzerr} contain the final
        deterministic-shrinkage quadrature summary.
  \item \code{success} and \code{termination_reason} must be read before
        treating the result as scientifically complete.
  \item \code{niter}, \code{nlive}, \code{n_likelihood_calls},
        \code{n_prior_calls}, and \code{n_proposals} describe cost.
  \item \code{config} stores the fully resolved immutable run configuration,
        including the resolved stopping policy.
  \item initial and final random-generator states and proposal metadata are
        retained as representations for auditability.
\end{itemize}

\subsection{Dead and live arrays}

The result separately stores points, \(\log L\), \(\log\pi\), \(\log q_0\),
\(\log\Psi_0\), and tie breakers for dead and final-live populations.  Dead
points additionally store \(\log X_i\) and their log quadrature contributions.
All arrays are copied, read-only, and validated.

\code{result.log_posterior_weights} has length \(K+N\) and follows the
dead-then-live ordering.  Convenience properties use the same order:

\begin{lstlisting}
points = result.all_points
log_psi0 = result.all_log_psi0
weights = result.posterior_weights
\end{lstlisting}

\subsection{Per-iteration history}

Every history field has length \code{niter}.  The history includes:

\begin{itemize}
  \item discarded threshold, \(\log X_i\), and \(\log\Delta X_i\);
  \item dead, live, and total \(\log Z\);
  \item information and theoretical \code{logzerr};
  \item remaining fraction, remaining log-evidence increment, live ESS,
        live-mean RSE, live log-evidence error, stability, and stopping streak;
  \item minimum, median, and maximum live \(\log\Psi\);
  \item proposal rows used for each completed replacement, cumulative
        likelihood calls, acceptance fraction, and elapsed seconds.
\end{itemize}

If no replacement completes, all history arrays are empty.  Diagnostic summary
helpers report NaN for unavailable final floating-point metrics and zero for
the final streak.

\subsection{Progress semantics}

\code{progress=True} displays an indeterminate live iteration count rather
than a percentage of the hard iteration limit.  It always shows live count,
likelihood calls, efficiency, current evidence, theoretical error, and stopping
streak.  Criterion-specific diagnostics appear only when enabled, proposal
status appears only after adaptive updates, and remaining fraction is omitted
from the terminal display.  A callable still receives the complete mapping
after every completed iteration, including a \code{criterion_<name>_met}
integer flag for every enabled criterion.

The acceptance fraction stored in history is
\[
  \frac{\text{completed replacements}}
       {\text{evaluated proposal rows}},
\]
where the denominator includes constrained-replacement proposal rows but
excludes the initial \(N\) live draws.  It is not the fraction of Python batch
calls and not the number of mathematically valid rows divided by proposals.

\section{Practical tuning and validation}

\subsection{Choosing the live count}

Increasing \(N\) makes the volume path in Equation~\eqref{eq:volume} finer,
increases the cost of initialization and every live-set diagnostic, raises the
maximum possible live ESS, and reduces the formal
\(\sqrt{\widehat H/N}\) scale for fixed \(H\).  It does not cure missing
proposal support.  Choose \(N\) by repeated-run calibration on targets similar
to the intended analysis.

\subsection{Choosing a proposal batch size}

A larger batch can improve vectorized likelihood throughput, but every row is
evaluated and counted even if the first row is valid.  Very large batches can
therefore waste expensive evaluations when acceptance is high.  Very small
batches can create Python/backend overhead.  The batch size does not change
the formal constraint or quadrature, but it can change the realized random
stream because unused rows from an accepted batch are not retained for the
next iteration.

\subsection{Recognizing rejection difficulty}

As the threshold rises, the probability under \(q\) of satisfying the
constraint may become small.  Warning signs are:

\begin{itemize}
  \item falling \code{history.acceptance_fraction};
  \item large \code{history.proposals};
  \item rapidly increasing likelihood calls per iteration;
  \item \code{constrained_sampling_exhausted}; or
  \item \code{plateau_stall} under strict ties.
\end{itemize}

Raising \code{max_proposals_per_replacement} only gives valid rejection more
time; it does not improve \(q\).  Proposal structure, bandwidth, posterior
coverage, parameterization, and plateau policy should be investigated.

\subsection{Recommended validation for a new setup}

\begin{enumerate}
  \item Verify the normalized prior analytically or numerically.
  \item Confirm the posterior sample column order and units against
        \code{parameter_names}.
  \item Inspect Morph metadata for selected groups and single parameters.
  \item Test \code{proposal.sample} and \code{proposal.log_prob} together in
        tail and boundary regions.
  \item Run multiple independent seeds and compare \(\log Z\), cost,
        termination reasons, and posterior summaries.
        To assess Morph-fitting variation as well as sampler variation, repeat
        the training/refitting stage rather than changing only the run seed.
  \item Where possible, compare with an analytic target, direct importance
        sampling under the same fixed \(q\), or another evidence estimator.
  \item Benchmark stopping policies on the target family rather than assuming
        the example hybrid thresholds are universal.
  \item Retain weighted samples and inspect posterior ESS and weight
        concentration before equal-weight resampling.
\end{enumerate}

\section{Implementation map}

\begin{longtable}{>{\raggedright\arraybackslash}p{0.32\linewidth}
                  >{\raggedright\arraybackslash}p{0.60\linewidth}}
\toprule
\textbf{Source file} & \textbf{Responsibility}\\
\midrule
\endfirsthead
\toprule
\textbf{Source file} & \textbf{Responsibility}\\
\midrule
\endhead
\code{src/nismo/model.py} &
Model protocol, callable adaptation, point-shape validation.\\
\breakcode{src/nismo/proposals/base.py} &
Normalized proposal sampling/density protocol.\\
\breakcode{src/nismo/proposals/morph.py} &
\MorphZ{} fitting, grouping inputs, fixed \code{GroupKDE} sampling and
row-wise normalized log-density adapter.\\
\code{src/nismo/constrained.py} &
Batch density evaluation, count tracking, support checks, and independent
strict/lexicographic constrained rejection.\\
\code{src/nismo/quadrature.py} &
Log-space deterministic volumes, dead/live contributions, evidence,
information, posterior weights, and theoretical error.\\
\code{src/nismo/stopping.py} &
Stopping configuration, finite-live diagnostics, criterion comparisons,
combination, iteration gate, and consecutive streak.\\
\code{src/nismo/config.py} &
Immutable run configuration, legacy stopping resolution, and validation.\\
\code{src/nismo/sampler.py} &
Serial state machine, hard-limit ordering, history/progress snapshots,
termination reason, finalization, and result construction.\\
\code{src/nismo/results.py} &
Immutable result/history containers, consistency checks, convenience arrays,
and systematic equal-weight resampling.\\
\code{src/nismo/diagnostics.py} &
Read-only summary diagnostics computed from a completed result.\\
\code{src/nismo/progress.py} &
Silent, callback, and optional \code{tqdm} progress reporters.\\
\bottomrule
\end{longtable}

\section{Concise user decision guide}

\begin{enumerate}
  \item \textbf{Do the training samples cover all posterior modes?}  If not,
        do not trust Phase~2 evidence; regenerate the posterior input.
  \item \textbf{Are prior and proposal densities normalized in the exact sample
        coordinates?}  If not, fix them before running.
  \item \textbf{Is the transformed pseudo-likelihood continuous?}  Use
        \code{tie_policy="strict"}.  If exact plateaus have nonzero
        probability, use \code{"randomized_plateau"}.
  \item \textbf{Is the posterior correlated?}  Select a suitable Morph grouping
        workflow rather than relying automatically on independent marginals.
  \item \textbf{Which scientific stop is intended?}  Use only \code{dlogz} for
        the remaining log-evidence increment, or only \code{stopping} for an
        explicit policy.
  \item \textbf{Did the run stop successfully?}  Require
        \code{result.success} before treating it as scientifically converged,
        and inspect the exact reason and final diagnostics.
  \item \textbf{Does the downstream tool accept weights?}  Use the native
        weighted representation if possible; resample only when necessary.
\end{enumerate}

\section{Summary of guarantees and non-guarantees}

Under the stated contracts, the current code provides:

\begin{itemize}
  \item a fixed normalized proposal used consistently for sampling and
        density evaluation;
  \item a clear change-of-measure identity;
  \item independent generation-order constrained rejection;
  \item deterministic log-volume quadrature with dead and live corrections;
  \item normalized posterior weights and reconstructible evidence;
  \item reproducible explicit RNG ownership;
  \item exact policy comparison, persistence, and termination status; and
  \item valid partial quadrature output for defined hard stops.
\end{itemize}

It does not guarantee:

\begin{itemize}
  \item posterior discovery or recovery of modes absent from training data;
  \item defensive proposal support;
  \item adaptation of \(q\) during the run;
  \item a complete uncertainty estimate in \code{logzerr},
        \code{live_logz_error}, or stability;
  \item efficient rejection at every threshold;
  \item that example stopping tolerances are calibrated for a new problem; or
  \item scientific convergence when a resource limit is reached.
\end{itemize}

\important{\textbf{The minimum safe interpretation rule is:}
check \code{result.success}, verify the exact \code{termination_reason},
inspect the final stopping and weight diagnostics, and validate the entire
setup over repeated independent runs.}

\end{document}
